% !TEX program = pdflatex
\documentclass[11pt,a4paper]{article}
\usepackage[T1]{fontenc}
\usepackage[utf8]{inputenc}
\usepackage{lmodern}
\usepackage[margin=25mm]{geometry}
\usepackage{amsmath,amssymb}
\usepackage{graphicx,booktabs,tabularx}
\usepackage[section]{placeins}
\usepackage[numbers,sort&compress]{natbib}
\usepackage{xurl}
\usepackage[hidelinks]{hyperref}
\usepackage{microtype}
\setlength{\emergencystretch}{2em}
\newcommand{\draftnotice}[1]{\par\smallskip\noindent\textit{Draft: #1}\par\smallskip}
\newcommand{\aicontribution}[5]{\par\nobreak\smallskip\noindent
\begin{minipage}{\linewidth}\small
\noindent\textbf{Question:} #1\par
\noindent\textbf{AI contribution:} #2\par
\noindent\textbf{Investigator choices:} #3\par
\noindent\textbf{Verification:} #4\par
\noindent\textbf{Evidence status:} #5\par\end{minipage}\par\smallskip}
\IfFileExists{generated/numbers.tex}{\input{generated/numbers.tex}}{}
\IfFileExists{generated/published_numbers.tex}{\input{generated/published_numbers.tex}}{}

\title{AI-assisted glioma research: a perspective with reproducible examples
in imaging and clinical evidence}
\author{Arvid Lundervold\\
\small Department of Biomedicine, University of Bergen, Norway}
\date{6 October 2026\\\small Working perspective draft}

\begin{document}
\maketitle

\begin{abstract}
Glioma research must connect anatomical tumor burden with the cells and neural
circuits that occupy the same brain. This perspective identifies two AI roles:
trained imaging models for compartment measurement, and generative research
assistance for constructing inspectable code, controls, figures and prose.
Our published MRI segmentation and subject-specific anatomical-profiling study
provides attributed imaging results and public-data illustrations. A retrospective
potassium-clearance case connects the investigator's question, assisted
implementation, independent checks and unresolved biological interpretation.
Reproducible CPU examples address imaging sensitivity, niche-dependent state
composition, potassium clearance, extracellular potential gradients and
tumor-growth feedback. Inputs and coefficients are prescribed, with explicit
units and separate timescales; none is a calibrated patient model. Experimental
literature motivates questions about astrocytes, neuron--glioma interactions
and cognitive function. Analog cognition and cytoelectric coupling remain
research perspectives, with field effects, cytoskeletal evidence and clinical
associations distinguished. A source-verified CANTAB case anchors cognition
in observed task performance. A proposed study asks whether preoperative
network/dynamics information adds to baseline cognition and anatomical burden.
Shared Python code generates notebook results, manuscript numbers and talk
figures. Conventional solvers reproduce specified equations without a language
model call. A proposed matched-task pilot separates quality-qualified completion
time, independently reviewed scientific quality and artifact relevance.
Biological calibration, human scientific review, clinical evaluation and
comparative measurement of AI benefit remain further work.
\end{abstract}

\noindent\textbf{Keywords:} glioma; tumor microenvironment; cancer neuroscience;
astrocytes; cognition; MRI; AI assistance; reproducibility

\section{Introduction}

The value of a computational glioma figure depends on the question it answers.
A segmentation may define a region for measurement, an equation may construct
a trajectory, and a clinical table may identify a candidate analysis population.
Each can be useful without supporting a treatment conclusion.

An image-analysis model, such as one developed using nnU-Net, maps an image
to labels and requires validation of that mapping~\cite{isensee2021}.
A generative research assistant can help formulate questions, write code,
inspect artifacts, and revise prose. A numerical solver advances a specified
equation. The solver is not an imaging model, and correctness of its update
does not validate the assistant that helped implement it.

Our focus is understanding glioma through AI-assisted imaging and modeling.
This perspective accompanies a talk entitled \emph{AI-Assisted Glioma
Research: Imaging, Tumor Dynamics and Clinical Evidence}. We organize the
examples by clinical research questions rather than by software products.
The aim is to connect a contribution to a check and a claim another
researcher can assess.

The central question is what happens at the interface between glioma and
surrounding brain. Imaging, tissue assays, and physiological measurements
observe different aspects of that interface. Our worked models ask how local
signals could change state composition or growth, and which observations
would be needed to distinguish these explanations. Evidence from these worked
models is synthetic; experimental findings are cited separately from our simulations.
A separately verified behavioral case supplies empirical cognitive observations.

The clinical motivation includes preservation of language and other cognitive
functions alongside control of tumor burden. Anatomical proximity, physiological
recruitment and task performance are distinct observations. Their integration
is a research question, rather than a consequence of segmentation accuracy.

\subsection{Two AI roles in a reproducible research workflow}
\textbf{Imaging AI} maps MRI to compartment labels. Its evidence concerns
the trained model's agreement with appropriate reference segmentations and
the geometry and meaning of its outputs. Our published SegResNet study
provides this component~\cite{lundervold2026}; its attributed results are
distinct from the synthetic measurement exercise.

\textbf{Generative research assistance} helps translate an investigator's
question into code, comparisons, figures and evidence-bounded prose. In this
project Codex contributed to these artifacts, as documented retrospectively.
The potassium example makes that contribution concrete without claiming an
independent biological discovery or an improvement over unassisted work.
The investigator supplies the research direction and remains responsible
for reviewing assumptions and interpretation. Development choices are not
presented as investigator-approved physiological estimates.

Conventional numerical solvers then reproduce the specified equations on
CPU. Rebuilding makes no language-model call. Each executable example below
identifies its question, AI contribution, investigator choices, verification
and evidence status. A check of the calculation supports implementation;
biological calibration and a comparative evaluation of assistance are
separate requirements. The proposed pilot in Section~\ref{sec:ai-pilot}
addresses research productivity and artifact quality, not patient outcomes.

\section{Measuring and locating a tumor}

\subsection{A published imaging foundation}

\aicontribution
{Where are MRI-defined compartments relative to subject anatomy (notebook 09)?}
{Trained SegResNet segmentation in the published study; research assistance here with attributed evidence and figure reuse.}
{The investigator requested this bridge; study methods belong to the source and interpretation needs review.}
{Published reference-mask Dice/HD95, compartment definitions, attribution and asset hashes.}
{Reported imaging study; no new inference, accuracy benchmark or clinical-utility result.}


Our recent feasibility study combines multiparametric MRI tumor segmentation
with subject-specific anatomical parcellation~\cite{lundervold2026}.
Intersecting compartments with anatomical regions produces volumes and
regional-burden \textit{hit-plots}. The companion code is available at
\url{https://github.com/arvidl/high-precision-glioma-segmentation}.
\IfFileExists{generated/published_numbers.tex}{
In the published UCSF-PDGM cohort of \PublishedImagingN{} subjects, reported
median Dice values were \PublishedDiceWT{} for whole tumor,
\PublishedDiceTC{} for tumor core, and \PublishedDiceET{} for enhancing tumor.
}{\draftnotice{Run make paper to load the attributed publication values.}}
These values describe agreement with the study's reference masks.
They are quoted published aggregates; that pipeline is not rerun here.
A longitudinal public-data illustration demonstrates descriptive change in
compartment volumes and regional burden.

The study's stated scope is feasibility and infrastructure, with no
patient-outcome or clinical-utility evaluation~\cite{lundervold2026}.
Its regional burden supplies anatomical context for biological questions,
without measuring malignant-cell programs, neuron--glioma synapses, or
functional eloquence. Relating those profiles to tissue assays, physiology,
or outcomes is proposed follow-up work. The executable construction below
teaches measurement sensitivity without importing patient images or labels.

\subsection{Methods, reported results and illustrations}
The published study uses co-registered T1, contrast-enhanced T1, T2 and FLAIR,
a trained deep-learning segmentation model (the unified MONAI BraTS
three-compartment SegResNet pipeline), and subject-specific
FreeSurfer anatomical parcellation~\cite{lundervold2026}. Native-space
intersections preserve each subject's anatomical labels. The study separates
segmentation agreement, anatomical localization and quantitative regional
profiling. Dataset reference masks provide the segmentation comparison;
they are not histological ground truth.

\begin{table}[htbp]
\centering
\caption{Reported median agreement in the published n=50 UCSF-PDGM
IDH-wildtype, WHO grade 4 glioblastoma cohort~\cite{lundervold2026}.
Values are the rounded abstract values, exported from a checked publication
record. No imaging inference or new benchmark was performed here.}
\begin{tabular}{@{}lrr@{}}\toprule
Compartment & Dice & HD95 (mm)\\\midrule
Whole tumor (WT) & \PublishedDiceWT & \PublishedHDWT\\
Tumor core (TC) & \PublishedDiceTC & \PublishedHDTC\\
Enhancing tumor (ET) & \PublishedDiceET & \PublishedHDET\\\bottomrule
\end{tabular}
\end{table}
WT combines NCR, ED/FLAIR abnormality and ET; TC combines NCR and ET.
The longitudinal LUMIERE illustration is one subject with six visits.
These results support feasibility of the measurement infrastructure, with
clinical outcomes and utility outside the study's evaluated scope.

\IfFileExists{generated/hpgs_mri_pipeline.png}{
\begin{figure}[htbp]
\centering\includegraphics[width=\textwidth]{generated/hpgs_mri_pipeline.png}
\caption{Published imaging illustration: multiparametric MRI, dataset-reference
compartment masks and subject-specific anatomy for public UCSF-PDGM-0085. Original
Figure 1 asset from Lundervold et al.~\cite{lundervold2026}, CC BY 4.0;
reused without altering the image. This illustration supplies anatomical
context, not a physiological or molecular measurement.}
\label{fig:published-pipeline}
\end{figure}
}{\draftnotice{Published illustrations need the explicit local
\texttt{make published-figures HPGS\_REPO=...} import from the author checkout.}}

\IfFileExists{generated/hpgs_regional_anatomy.png}{
\begin{figure}[htbp]
\centering\includegraphics[width=\textwidth]{generated/hpgs_regional_anatomy.png}
\caption{Published public UCSF-PDGM-0020 illustration of tumor contours
and native anatomical labels. Original Figure 4a--b asset from
Lundervold et al.~\cite{lundervold2026}, CC BY 4.0. Anatomical labels
locate a possible functional question; they do not measure eloquence.}
\label{fig:published-anatomy}
\end{figure}
}{}
\IfFileExists{generated/hpgs_hitplot_reference.png}{
\begin{figure}[htbp]
\centering\includegraphics[width=0.94\textwidth]{generated/hpgs_hitplot_reference.png}
\caption{Published regional overlap volumes for public UCSF-PDGM-0020,
original Figure 4c asset~\cite{lundervold2026}, CC BY 4.0. NCR, ED and ET
retain the source legend and mL units. The upstream label \texttt{source=gt}
identifies dataset reference masks, rather than deep-learning predictions
or histological ground truth. The displayed top twelve regions are the
published selection, not a new analysis or an outcome association.}
\label{fig:published-hitplot}
\end{figure}
}{}
The image assets remain local and are checked against source hashes.
Their author-manuscript revision is recorded separately from the cleaned
public code revision, which does not distribute these manuscript assets.
No restricted BGO-D illustration or acquired patient file is imported.

\subsection{Constructed compartments}

\aicontribution
{How do geometry and mask perturbations change measurement (notebook 01)?}
{Research assistance adapting unit-explicit measurement code and teaching controls.}
{Imaging was requested; grid, labels and perturbations are prescribed construction choices.}
{Affine volume, Dice, matched grids and refusal of ambiguous geometry.}
{Synthetic masks; independent of the trained imaging model.}


The imaging construction is an ellipsoid on a
$1.2\times1.2\times2.0$\,mm grid. Nested compartments use explicit labels
1, 2, and 4 as constructed analogues of necrosis or non-enhancing tumor,
edema or FLAIR abnormality, and enhancing tumor, arranged as in a typical
glioblastoma: necrotic core, enhancing rim, surrounding edema. These definitions describe
a construction; they are not histologically observed tissue.

For a binary region $M$ and a voxel-to-world affine $A$ in millimetres,
\begin{equation}
 V(M)=\frac{|M|\,|\det A_{1:3,1:3}|}{1000}\quad\mathrm{mL}.
\end{equation}
The determinant handles rotation, reflection, and shear. The NIfTI reader
requires a coded transform, established units, and finite integer labels.
It rejects conflicting coded qform and sform transforms. Comparisons require
the same physical grid; the program does not silently register or resample.

A four-voxel translation is compared with a 12\% increase in each ellipsoid
radius. These perturbations were prescribed. Translation changes overlap
without changing voxel count in this construction. Expansion changes the
selected volume.

\IfFileExists{generated/numbers.tex}{
The reference volume is \ExampleReferenceML{}\,mL.
The translated mask has Dice \ExampleTranslatedDice{}, and the expanded
mask has Dice \ExampleExpandedDice{} with a volume increase of
\ExampleExpandedPercent{}\%.
}{\draftnotice{Run make paper to generate the measurement values.}}

\IfFileExists{generated/imaging.pdf}{
\begin{figure}[htbp]
\centering\includegraphics[width=\textwidth]{generated/imaging.pdf}
\caption{Overlap and volume respond differently to two small segmentation
errors (synthetic construction). \textbf{(A)} Mid-plane of a constructed
ellipsoid (semi-axes 10, 12 and 16\,mm) on a $64\times64\times40$ grid of
$1.2\times1.2\times2.0$\,mm voxels, with three nested, mutually exclusive
labels in BraTS numbering: label~1, necrotic core or non-enhancing tumor
(purple); label~4, enhancing rim (orange); label~2, edema or FLAIR
abnormality (teal). The labels are prescribed regions, not observed tissue
or the output of a trained network. \textbf{(B)} Reference whole tumor
(labels 1+2+4; grey; \ExampleReferenceML{}\,mL) with the same mask
translated by four voxels (4.8\,mm) along $x$ (magenta) and an ellipsoid
with every semi-axis enlarged by 12\% (blue; about one voxel at the
boundary). \textbf{(C)} Whole-tumor volume, signed volume error $\Delta V$
and Dice relative to the reference. The translation preserves volume but
lowers Dice to \ExampleTranslatedDice{}; the expansion keeps Dice at
\ExampleExpandedDice{} yet increases volume by \ExampleExpandedPercent{}\%,
close to the continuous expectation $1.12^3-1\approx40.5\%$. Volumes use
voxel counts and the affine determinant; Dice is computed on the identical
grid without registration or resampling. Axes give millimetres from the
ellipsoid centre, not patient anatomy.}
\label{fig:imaging}
\end{figure}
}{\draftnotice{Synthetic imaging figure has not been generated.}}

\subsection{Clinical relevance and proposed public example}

An overlay should be interrogated before interpreting a number: which
compartment, label convention, grid, and source produced it? Anatomical
context can inform a research question, while structural proximity alone
does not establish functional risk. Generated inpainted anatomy must remain
identified as generated.

The reused published public-data illustrations demonstrate the study's
anatomical profiling. A new case analysis would still require source images,
preparation records and independent imaging review. The reuse does not
constitute a new review of segmentation or anatomical accuracy.

\section{Cells, neural dynamics and cognition at the glioma--brain interface}

\subsection{Biological questions and levels of evidence}

The tumor microenvironment includes immune and vascular compartments and
the surrounding neural tissue. Macrophage-associated signaling can drive
MES-like glioblastoma cell-state transitions in experimental systems
~\cite{hara2021}; the associated figure correction is recorded
~\cite{hara2025correction}. Hypoxia-associated macrophage niches and
vascular interactions illustrate the spatial nature of this ecosystem
~\cite{wang2024}. These findings motivate biological questions rather
than supplying coefficients for our constructions.

Cancer neuroscience offers a bridge to neurosurgery and tumor biology
~\cite{winkler2023,monje2025}.
Synaptic and electrical neuron--glioma interactions have experimental
support~\cite{venkatesh2019}. Neuron--glioma synapses show plasticity
regulated by BDNF--TrkB signaling, and greater depolarizing current promotes
glioma proliferation in experimental models~\cite{taylor2023}. Human
cortical-tissue experiments connect
neuronal excitability, glioma synaptic activity, and proliferation
~\cite{mcalpine2026}. A study combining recordings, biopsies, and experimental
work also links tumor functional connectivity with cognition and patient
survival~\cite{krishna2023}. Its patient associations do not demonstrate that
altering connectivity improves survival. Evidence from particular glioma
subtypes and preparations should retain those boundaries.

Tumor microtubes add a cellular network level~\cite{osswald2015}. In
experimental models they help glioma networks resist surgical lesions and
chemotherapy~\cite{weil2017}, carry neuron--glioma synapses
~\cite{venkataramani2019}, take part in neuron-like brain invasion
~\cite{venkataramani2022}, and connect rhythmically active pacemaker-like cells
whose KCa3.1-dependent activity drives network growth~\cite{hausmann2023}.
They are detected mainly by intravital two-photon microscopy in mouse models,
with confirmation in resected tissue, and are far below MRI resolution. A
transcriptional connectivity signature that includes CHI3L1 now estimates
network connectivity in patient samples; high connectivity is associated with
the mesenchymal subtype and shorter survival~\cite{hai2024}.

\subsection{Astrocytes as physiological and tumor partners}
Non-malignant astrocytes regulate the extracellular environment and support
synaptic activity. Hippocampal-slice experiments with glial-conditional Kir4.1
perturbation link astroglial potassium clearance to short-term synaptic
responses~\cite{sibille2014}. These physiological observations motivate a
clearance example; they do not supply its assumed rates or demonstrate the
same behavior in tumor-infiltrated human cortex.

Astrocyte--glioblastoma connections also have metabolic consequences.
Watson and colleagues demonstrated GAP43-dependent astrocyte-to-glioblastoma
mitochondrial transfer and enhanced tumorigenicity in experimental systems
~\cite{watson2023}. The work includes University of Bergen investigators.
That process is distinct from potassium clearance. An astrocyte-like
malignant-cell transcriptional program likewise does not identify a
non-malignant or reactive astrocyte. Cell identity, state and compartment
must remain explicit when proposing a model.

\subsection{Local field potentials and electric fields}
LFPs record extracellular potential relative to a reference electrode.
Transmembrane currents, especially synaptic currents, together with cellular
geometry, synchrony and volume conduction contribute to the measured signal
~\cite{buzsaki2012}. A potential in volts and its spatial gradient in volts
per metre are different quantities:
\begin{equation} \mathbf E=-\nabla\phi.\label{eq:field}\end{equation}
A temporal trace at one electrode does not determine that gradient.
Cortical-slice experiments establish that weak extracellular fields can
influence membrane potential and spike timing under the tested conditions
~\cite{anastassiou2011}. These ephaptic effects do not alone establish
cortex-wide cognitive organization or a field-mediated glioma mechanism.

\subsection{An empirical clinical anchor: source-verified CANTAB}
\label{sec:cantab}

\aicontribution
{What cognitive performance was observed before surgery and at follow-up (notebook 10)?}
{Research assistance retrieving and pinning the behavioral release, checking score precision, and implementing the shared figure.}
{The investigator requested an empirical anchor; four cognitive domains were selected for teaching, with interpretation pending review.}
{Source SHA-256, explicit case match, visit chronology, units, missingness, diagnosis groups and agreement at released precision.}
{Descriptive real-data example; no reliable-change classification, causal surgical effect or evaluated predictor.}

The CANTAB behavioral release supplies an empirical anchor distinct from our
synthetic neuroglial models. Aerts and colleagues assessed patients and caregivers
before surgery and several months afterwards~\cite{aerts2019,aertsosf}.
We pinned the released data, legend and non-imputed analysis code, preserving
source definitions and day/month/year assessment dates. The behavioral resource
is distinct from the selected T1, diffusion and resting-state BOLD MRI
cohort~\cite{aerts2022}; participant linkage and imaging QC cannot be inferred
from equal row positions or assumed cohort identity.

\IfFileExists{generated/cantab_numbers.tex}{\input{generated/cantab_numbers.tex}}{}
\ifdefined\CantabBehavioralN
The behavioral table contains \CantabBehavioralN{} records: \CantabGliomaN{}
glioma, \CantabMeningiomaN{} meningioma and \CantabControlN{} caregiver controls.
Complete paired scores across the four cognitive domains are available for
\CantabGliomaPairedN{}, \CantabMeningiomaPairedN{} and \CantabControlPairedN{}
participants, respectively. These are code-generated behavioral availability
counts, before any imaging linkage or QC exclusions. Diagnoses remain separate.
\fi

The illustration uses RTI five-choice reaction time in milliseconds, RVP
sensitivity (A prime), SOC problems solved in minimum moves, and SSP spatial
span. They retain their own scales. Lower reaction time means a faster
response; higher sensitivity, planning count and span generally indicate better
task performance. Neither a directional difference nor a percentage change
establishes clinically reliable change. Compatible normative and retest
information is not supplied by this release. The small repeated caregiver
sample supports exploratory comparison, rather than a validated normative
database or isolation of a surgical effect.

\IfFileExists{generated/cantab_case.pdf}{
The investigator's local case table was checked against one explicitly matched
glioma row. \CantabVerifiedMetricN{} of its \CantabLocalMetricN{} CANTAB measures
are represented in this release, including motor screening.
Of the paired shared cells, \CantabExactCells{} agree exactly and
\CantabPrecisionCells{} agrees at released decimal precision. The figure uses
released values throughout. In particular, the small attention difference in
the more precise local table is not resolved by the release; the displayed
attention scores are equal. The assessments were \CantabCaseInterval{} days
apart. This is an assessment interval, not a measured postoperative interval.
\begin{figure}[htbp]
\centering\includegraphics[width=\textwidth]{generated/cantab_case.pdf}
\caption{Source-verified single-case CANTAB example. Each panel retains its
native scale and the released decimal precision; values come from code, not
manual transcription. Assessment interval: \CantabCaseInterval{} days.
RTI uses the released five-choice measure without guessing mean versus median.
The paired scores describe observations across surgery, recovery and intervening
care. Reliable change and causal effects are not assessed. Source:
OSF Data.csv revision 2 and Legend.csv revision 1~\cite{aertsosf}, CC BY 4.0.
The local source-case identifier and exact dates are omitted.}
\label{fig:cantab}
\end{figure}
}{\draftnotice{The empirical CANTAB figure requires explicit acquisition and
local case configuration, followed by make cantab. No synthetic case substitutes
for missing empirical inputs.}}

Treatment, seizures, medication, fatigue, emotional state, practice effects
and assessment conditions may contribute to the observed differences.
Two perioperative visits cannot identify tumor-growth effects or a pure
resection effect. A prospective growth investigation needs serial imaging and
cognition with appropriate treatment timing. This example contains no
electrophysiological, astrocytic or microtubule measurements.

Our published segmentation methods provide an imaging and quality-control
bridge~\cite{lundervold2026}. Their T1/T1c/T2/FLAIR SegResNet inputs are not all
available in BTC T1/diffusion/BOLD. The four-input trained model therefore
cannot be directly applied to that resource. No missing contrast is synthesized
or silently substituted; supplied masks, T1 anatomy and connectomics require
their own reviewed provenance.

\subsection{Cognitive function and the analog-cognition perspective}
The language-task recordings and tumor studies of Krishna and colleagues
provide a disease-specific connection among glioma, circuit remodeling and
cognitive performance~\cite{krishna2023}. Their patient language and survival
associations retain the limitations of association, including causal direction.
Cognition therefore belongs in the clinical research agenda as an explicitly
assessed outcome, rather than a label attached to a simulated activity index.

Miller, Brincat and Roy propose a hybrid synapse/spike/wave framework in which
rhythmic fields and wave dynamics organize context-dependent neural activity
~\cite{miller2026}. The article is a theoretical review. It motivates asking
how glioma-related perturbations of organized neural dynamics might affect
cognition beyond anatomical burden; that glioma application is our proposed
question. Its broader consciousness interpretation remains a theoretical
horizon, with language and other prespecified cognitive outcomes as the
practical focus here. The investigator's September 2026 HBF presentation is
an educational companion, rather than an independent biological validation.

\subsection{Cytoelectric coupling and cytoskeletal evidence}
Pinotsis, Fridman and Miller propose that electric fields could influence
cytoskeletal organization through electrodiffusion and mechanotransduction
~\cite{pinotsis2023}. Their cytoelectric-coupling perspective explicitly
calls for LFP recordings linked to molecular-level imaging as a direct test.
In prefrontal recordings, Pinotsis and Miller report that field-to-neuron
(ephaptic) influences were stronger than neuron-to-field influences and
correlated with trial-by-trial changes in oscillatory power~\cite{pinotsis2026};
the finding concerns neural activity, not glioma cells.
A field-to-microtubule-to-glioma-growth or cognition pathway is therefore
a proposed hypothesis in this manuscript, without a fitted coupling term.

A separate glioblastoma study disrupted cytoplasmic connexin-43--microtubule
interactions and reduced glioblastoma stem-like cell survival and xenograft
growth~\cite{smyth2025}. It supplies experimental cytoskeletal relevance;
it does not establish a physiological LFP-driven mechanism. Intracellular
microtubules are tubulin polymers, whereas tumor microtubes are cell
projections supporting intercellular glioma networks~\cite{osswald2015}.
Their names must not be used interchangeably.

\section{Reproducible constructions across timescales}
The examples deliberately resolve different quantities and timescales.
They are not parts of one calibrated multiscale simulator. In particular,
seconds-scale extracellular potassium, a spatial potential illustration
and a day-scale dimensionless growth signal are not interchangeable.

\subsection{A niche signal changes composition at fixed abundance}

\aicontribution
{Can composition change at fixed local abundance (notebook 05)?}
{Research assistance implementing the model, comparisons, tests and figures.}
{Biological modeling was requested; states, reversibility, cue and rates await scientific review.}
{Exact solution, conservation, equilibrium, no-cue limit and partition invariance.}
{Synthetic composition; no observed phenotype or fitted niche effect.}


Let $P$ aggregate other tumor states and $M$ represent a MES-like program
analogue, both dimensionless normalized abundance. A prescribed spatial
cue $s\in[0,1]$ gives
\begin{equation}
 \frac{dM}{dt}=k_f sP-k_rM,\qquad
 \frac{dP}{dt}=-\frac{dM}{dt}.
 \label{eq:niche}
\end{equation}
The model conserves $P+M$ locally. The cue is neither an immune-cell count
nor a measured cytokine concentration. Reversibility is an assumed teaching
choice; no transition rate is inferred from the cited experiments.
The fixed-cue system has an exact solution, used by the package.
Growth, migration, death, and immune-cell dynamics are excluded.

A constructed 61-by-61 plane uses 1\,mm spacing, a Gaussian abundance profile,
an initially uniform 10\% MES-like fraction, and an off-center Gaussian cue.
With $k_f=0.08$/day and $k_r=0.02$/day over ten prescribed days,
\IfFileExists{generated/numbers.tex}{
the overall MES-like fraction becomes \ExampleNicheFractionPercent{}\%
with the cue and \ExampleNicheNoCuePercent{}\% in the no-cue comparator.
}{\draftnotice{Run make paper to generate the niche values.}}
These are constructed compositions, not patient measurements.
Both scenarios retain the same total abundance.

\IfFileExists{generated/microenvironment.pdf}{
\begin{figure}[htbp]
\centering\includegraphics[width=\textwidth]{generated/microenvironment.pdf}
\caption{Synthetic niche construction. The cue changes the distribution
of a MES-like program analogue while total normalized abundance is fixed.
The line plot compares a prescribed cue with its absence under the same
reverse rate. Axes, indices, and rates are constructed; neither molecular
state nor MRI appearance is inferred.}
\label{fig:niche}
\end{figure}
}{\draftnotice{The niche figure has not been generated.}}

This example separates abundance from state composition. It supplies no
MRI forward model and cannot establish that two biological tumors would
have identical scans. A molecular or protein assay would be required
to investigate the states; a suitable perturbation would be needed to
test a causal signaling explanation.

\subsection{Seconds-scale astrocytic potassium clearance}

\aicontribution
{What changes under fixed release and altered clearance (notebook 07)?}
{Codex helped formulate and implement the excess-K/Nernst example and checks.}
{The investigator requested astrocytes; volume, rates, pulses and concentrations are prescribed implementation choices.}
{Independent analytic solution, concentration balance, limiting cases, clearance ambiguity and Nernst limits.}
{Synthetic physiology; actual contribution summarized retrospectively, with scientific review pending.}

Notebook 07 introduces a one-compartment excess extracellular K concentration
$k$ in mM above a fixed 3\,mM reference:
\begin{equation}
 \frac{dk}{dt}=q(t)-(\lambda_a+\lambda_o)k.
 \label{eq:potassium}
\end{equation}
Release $q$ is prescribed in mM/s; lumped astrocytic and other clearance rates
are in s$^{-1}$. The exact constant-input solution advances each interval.
The receiving reservoir is implicit; Kir4.1 kinetics, spatial buffering,
intracellular storage, rerelease and firing are omitted. Cumulative release,
retained excess and clearance share the same fixed extracellular reference
volume, supporting an explicit flux-balance check.

Release is 0.6\,mM/s during $[2,6)$ and $[10,14)$ seconds and zero otherwise.
Other clearance is fixed at 0.1\,s$^{-1}$, while astrocytic clearance is
0.8, 0.2 or 0\,s$^{-1}$. Starting from zero excess, 0.01-second intervals
cover 24 seconds. These are prescribed conditions, not fitted healthy or
impaired astrocytes. The ideal K-selective Nernst equilibrium uses a fixed
intracellular concentration of 140\,mM and 37\,$^\circ$C:
\begin{equation}
 E_K=\frac{RT}{F}\ln\frac{[K]_{out}}{[K]_{in}}.
\end{equation}
Concentrations replace activities as an ideal-solution assumption.
$E_K$ is neither a simulated membrane voltage nor a cognitive endpoint.

\IfFileExists{generated/numbers.tex}{
Peak excess concentrations are \ExampleKReferencePeak{} and
\ExampleKReducedPeak{}\,mM for reference and reduced astrocytic clearance.
Their maximal equilibrium shifts are \ExampleKReferenceShift{} and
\ExampleKReducedShift{}\,mV.
}{}

\IfFileExists{generated/astrocyte_potassium.pdf}{
\begin{figure}[htbp]
\centering\includegraphics[width=\textwidth]{generated/astrocyte_potassium.pdf}
\caption{Synthetic potassium-clearance construction. Release intervals
are shaded and all scenarios retain identical release and other clearance.
The right panel computes a change in ideal K equilibrium voltage, not
membrane potential, firing or cognition. No glioma or astrocyte observation
was used to set the coefficients.}
\label{fig:potassium}
\end{figure}
}{}

The K trajectory depends only on $\lambda_a+\lambda_o$. Alternative
clearance attribution with the same sum gives the same trajectory, a
numerically checked ambiguity. Measuring K alone would not identify
astrocytic clearance within this model family. The construction is separate
from the slow activity-associated tumor-growth model below.

\subsubsection{Worked AI-assisted development case: retrospective evidence}
The saved development summary records the investigator's request to include
astrocytes, LFPs and cognition. Codex helped narrow that direction to the
clearance question and implement the package, comparison, tests and figures.
The one-compartment equation and rates are inspectable implementation choices,
not recorded approvals by the investigator. The original development record
is preserved; \path{docs/worked-ai-potassium.md} connects it to source and
verification. Notebook 04 provides the compact executable demonstration,
and notebook 07 retains the full comparison.

The chain is question, assisted implementation, checks, interpretation and
remaining uncertainty. In particular, accounting verifies
\begin{equation}
 k(0)+\int_0^t q(u)\,du
 =k(t)+C_a(t)+C_o(t),
\end{equation}
where both cleared quantities use the fixed extracellular reference volume.
An independent constant-input solution supplies a second check.
Equal-total-clearance ambiguity constrains interpretation even when both
checks pass. Physiological measurements and selective perturbations would
be needed to distinguish release from clearance contributions.

This retrospective case documents a contribution and its acceptance evidence;
it supplies no captured dialogue, rejected suggestion, elapsed-time comparison
or independent scientific approval. Human biological review remains pending,
and the Nernst calculation supplies no membrane-voltage or cognitive prediction.

\subsection{Extracellular potential as an observation problem}

\aicontribution
{Does potential identify a spatial field or cognitive mechanism (notebook 08)?}
{Research assistance with the spatial derivative, controlled examples and evidence map.}
{The investigator requested LFPs/cognition; profiles and sampling are prescribed teaching choices.}
{Analytic derivative, sign, V/m conversion, offset invariance and primary-source scope.}
{Synthetic static fields; cited experiments and hypotheses remain separate.}

Notebook 08 prescribes two static one-dimensional sinusoidal potentials on
an explicit $[-5,5]$\,mm line. Both have 100\,$\mu$V amplitude; wavelengths
are 10 and 2.5\,mm. A second-order finite-difference derivative converts
microvolts and millimetres to one field component in V/m. A common voltage
offset leaves that derivative unchanged. The analytic derivative supplies
an independent comparison; no neuronal current source or volume conductor
is inferred.

\IfFileExists{generated/numbers.tex}{
The analytic peak fields are \ExampleFieldLong{} and
\ExampleFieldShort{}\,V/m.
}{}

At the central position both potentials are zero, while their gradients
differ by approximately a factor of four. Equal potential amplitudes or a
single instantaneous voltage therefore do not identify the field gradient.
A real estimate would need spatial sampling, reference and noise treatment,
a defined conductivity/source model where relevant, and an account of
cellular geometry. The illustration has no temporal oscillation frequency.

\IfFileExists{generated/lfp_field.pdf}{
\begin{figure}[htbp]
\centering\includegraphics[width=\textwidth]{generated/lfp_field.pdf}
\caption{Synthetic spatial potentials and their gradients. These curves
are prescribed, not recorded LFPs or a cortical wave reconstruction. The
computed field is one spatial component; intracellular fields, ephaptic
responses, microtubule changes, tumor growth and cognition are not computed.}
\label{fig:lfp}
\end{figure}
}{}

This example clarifies a measurement requirement for any future field-based
hypothesis. It supplies no coefficient linking fields to the cytoskeleton
or to the normalized growth signal.

\subsection{Day-scale reaction--diffusion baseline}

Serial MRI observes a glioma through boundaries, delineated manually or
automatically and each with error. Between examinations they may expand,
shift, shrink after treatment or split into separate regions. Whether such
observations can support forecasts of growth in space and time has been
framed by analogy with numerical weather prediction~\cite{yankeelov2015}.
Weather forecasting now combines physics-based ensembles with machine-learned
models that match or exceed operational skill, including for tropical cyclone
tracks~\cite{price2025,bodnar2025,rasp2026}. Those models learned from decades
of dense, global observations assimilated by physics-based models. A glioma is
seen in a few scans months apart, its course is deliberately altered by
treatment, and pseudoprogression after radiotherapy can mimic
growth~\cite{wen2023}. Unlike a compact cyclone with a trackable centre, a
diffuse glioma's infiltrating margin is the part that is hardest to see.

The tumor is also coupled to the brain around it. Neuronal activity can
promote glioma growth through synaptic and electrical
integration~\cite{venkatesh2019}; tumor cells form microtube-connected
networks~\cite{osswald2015}; tumor functional connectivity is associated with
cognition and survival~\cite{krishna2023}; and compression injures neurons and
promotes glial neuroinflammation in experimental models~\cite{zarodniuk2026}.
A clinically useful forecast would concern this tumor--brain system, including
function and cognition, not only a boundary. The baseline below is a first,
deterministic step, with no data assimilation, ensemble, learned component or
coupling to the brain.

\subsubsection{A prescribed two-dimensional experiment}

\aicontribution
{How do diffusion and growth assumptions change density (notebook 02)?}
{Research assistance adapting solver interfaces, controls and explanation.}
{Tumor dynamics was requested; coefficients, geometry and boundaries are prescribed.}
{No-change/diffusivity comparators, conservation without growth, analytic limits and step constraints.}
{Synthetic day-scale model; no fitted patient trajectory.}


The baseline dynamics construction advances normalized density $c$ on a synthetic plane:
\begin{equation}
\frac{\partial c}{\partial t}
=\nabla\cdot(D\nabla c)+\rho c(1-c).
\label{eq:dynamics}
\end{equation}
This is the Fisher--Kolmogorov--Petrovsky--Piskunov equation, whose fronts
travel at speed $2\sqrt{D\rho}$~\cite{fisher1937,kpp1937}; Murray develops its
biological and brain-tumor applications~\cite{murray2003}. Swanson, Alvord and
Murray used the proliferation--invasion form with faster migration in white
than grey matter~\cite{swanson2000} and examined its implications for
treatment~\cite{swanson2004}. Our logistic term is a bounded teaching variant.
Diffusivity $D$ is in mm$^2$/day, proliferation $\rho$ is per day, and
coordinates are millimetres. A binary disk of radius 7\,mm is the initial
condition. The boundary is no flux. The explicit finite-difference update
uses $D=0.35$, $\rho=0.035$, a 1\,mm grid, and a 0.25-day step for 20 days.
The comparison doubles $D$ while retaining the other inputs.
This is the compact construction used for the figure; notebook 02 extends it
with an initial occupancy of 0.65, 28 days, separate reaction and diffusion
controls, and a comparison of $D$ and $\rho$ with published patient
estimates~\cite{wang2009}.

The implementation rejects time increments outside its diffusion stability
bound,
\begin{equation}
\Delta t\leq
\frac{1}{2\max(D)(\Delta x^{-2}+\Delta y^{-2})}.
\end{equation}
Tests include conservation of total density for $\rho=0$, invalid steps,
and directional diffusivity. Thresholds 0.5 and 0.05 are display choices;
neither is calibrated to MRI visibility. Coefficients are prescribed, not
estimated from a patient's scans.

\IfFileExists{generated/numbers.tex}{
The constructed initial area is \ExampleInitialArea{}\,mm$^2$.
The prescribed solution covers \ExampleVisibleArea{}\,mm$^2$ at threshold
0.5 and \ExampleLowArea{}\,mm$^2$ at threshold 0.05.
The doubled-diffusivity solution covers
\ExampleAlternativeArea{}\,mm$^2$ at threshold 0.5.
}{\draftnotice{Run make paper to generate the dynamics values.}}

\IfFileExists{generated/dynamics.pdf}{
\begin{figure}[htbp]
\centering\includegraphics[width=\textwidth]{generated/dynamics.pdf}
\caption{Constructed density and prescribed dynamics after 20 days.
The middle and right panels use different prescribed diffusivities.
Cyan and white contours indicate thresholds 0.05 and 0.5.
The initial disk is a no-change comparator. This is a numerical experiment,
not an MRI prediction or an inferred growth rate.}
\label{fig:dynamics}
\end{figure}
}{\draftnotice{Synthetic dynamics figure has not been generated.}}

\subsubsection{Connecting the experiment to longitudinal MRI}

\aicontribution
{How does a prescribed step compare with a later prepared label (notebook 03)?}
{Research assistance adapting the local recipe; trained labels are supplied upstream.}
{Prepared-case selection, anatomy and permitted inputs require investigator review.}
{Explicit preparation, geometry, labels, visit interval and numerical comparators.}
{Optional single-slice workflow, not executed here; no new response or recurrence result.}


LUMIERE supplies longitudinal glioblastoma MRI, automated segmentations,
and expert RANO evaluations~\cite{suter2022}. The optional local recipe
adapts the predecessor's Patient-048 slice workflow. It expects an explicitly
configured prepared grid and MONAI bundle label records. It verifies the
geometry and label scheme but performs no segmentation or registration.
Its step is between two later visits and does not cross the resection.

That recipe is available for reproduction; this initial draft does not
import the predecessor's real-data numbers as new results. A checked
six-visit montage and native-space compartment trajectory are proposed
extensions. They would require a documented case-selection rationale,
preparation records, and review of the source images.

Measured imaging change and clinical response require different evidence.
RANO 2.0 provides criteria for glioma trials that incorporate imaging
alongside clinical and treatment context~\cite{wen2023}. It remains the
current MRI-based framework. PET RANO 1.0 adds amino-acid PET response
criteria~\cite{albert2024}, updated RANO/EANO recommendations cover clinical
PET use in glioma~\cite{galldiks2025}, and a progress report describes the
wider RANO effort, including extent-of-resection classification and an AI
working group~\cite{vandenbent2025}.
Overlap with a later automated label does not establish recurrence
prediction, response category, or a mechanism of treatment resistance.

\subsection{A coarse-grained neuron--glioma feedback experiment}

\aicontribution
{What separates one-way growth modulation from feedback (notebook 06)?}
{Research assistance implementing the model, nested controls and sensitivity.}
{Cancer neuroscience was requested; the slow index and coefficients are prescribed choices.}
{Uncoupled limit, analytic constant-signal case, conservation, refinement and growth confounding.}
{Synthetic day-scale density/signal, with no fitted physiology or cognitive endpoint.}


The growth-focused construction adds a slow activity-associated growth-signal index
$a\in[0,1]$ to normalized density $c$:
\begin{align}
 \frac{\partial c}{\partial t}
 &=\nabla\cdot(D\nabla c)+\rho(1+\kappa a)c(1-c),\\
 \frac{\partial a}{\partial t}
 &=\frac{b-a}{\tau}+\beta c(1-a).
 \label{eq:neural}
\end{align}
Here $b$ is a prescribed dimensionless external drive, $\kappa$ is a
dimensionless signal-to-growth coupling, $\beta$ is a density-to-signal
feedback rate per day, and $\tau$ is a relaxation time in days.
The index $a$ is not voltage, firing frequency, EEG, or synapse count.
Its day-scale relaxation represents a constructed slow signal; the
model resolves neither electrophysiology nor specific growth factors.
These equations are our phenomenological teaching choice, not an
experimentally validated reduction of the cited biological studies.

The plane and spacing match the niche illustration, but the initial density
and drive are separately constructed Gaussians. Initial density peaks at
0.35, and $b$ ranges from approximately 0.1 to 0.8. We prescribe
$D=0.2$\,mm$^2$/day, $\rho=0.04$/day, $\kappa=1.5$, $\beta=0.12$/day,
and $\tau=4$\,days. A 0.1-day explicit step advances twenty prescribed days.
Density has no-flux boundaries; the signal has no transport.
Both updates use the previous fields.

Three nested comparisons hold the initial density, drive, diffusion,
basal growth, grid, and duration fixed:
\begin{enumerate}
 \item Baseline: $\kappa=\beta=0$, recovering the existing density solver.
 \item One-way: $\kappa=1.5$, $\beta=0$, retaining a fixed growth-modulating signal.
 \item Feedback: $\kappa=1.5$, $\beta=0.12$/day, allowing density to alter the signal.
\end{enumerate}
The explicit implementation rejects steps outside the sufficient bound
\begin{equation}
 \Delta t\leq\min\left\{
 [2D(\Delta x^{-2}+\Delta y^{-2})+\rho(1+\kappa)]^{-1},
 [\tau^{-1}+\beta]^{-1}\right\}.
\end{equation}
This includes reaction and feedback, not only diffusion. Inputs are bounded
and no output is silently clipped. Tests check the uncoupled limit,
conservation without reaction, an analytic constant-signal logistic solution,
and step refinement.

The integral $\sum c\,\Delta x\Delta y$ has units mm$^2$ on this plane.
It is not cell number or clinical tumor volume.
\IfFileExists{generated/numbers.tex}{
Its final values are \ExampleNeuralBaselineIntegral{},
\ExampleNeuralOneWayIntegral{}, and \ExampleNeuralFeedbackIntegral{}\,mm$^2$
for baseline, one-way, and feedback. The feedback-to-baseline ratio is
\ExampleNeuralFeedbackRatio{}. Halving the step changes the feedback
integral by \ExampleNeuralStepDifferencePercent{}\%.
The additional feedback changes the integral by
\ExampleNeuralFeedbackOverOneWayPercent{}\% relative to the one-way case,
so the total comparison with baseline must not be attributed entirely to
feedback.
}{\draftnotice{Run make paper to generate the neural-model values.}}
These differences describe the prescribed equations and inputs.
They do not estimate a neuronal effect in a patient.

\IfFileExists{generated/neuron_glioma.pdf}{
\begin{figure}[htbp]
\centering\includegraphics[width=0.88\textwidth]{generated/neuron_glioma.pdf}
\caption{Synthetic neuron--glioma construction. Final baseline and feedback
densities share a color scale. The signal-change panel compares the feedback
index with its prescribed initial drive. The traces include the one-way
comparator, separating growth modulation from the additional feedback
assumption. No field is a measured neural or imaging signal.}
\label{fig:neural}
\end{figure}
}{\draftnotice{The neuron--glioma figure has not been generated.}}

\subsection{Sensitivity and identifiable explanations}

A small grid varies $\kappa$ over 0, 0.75, and 1.5 and $\beta$ over
0, 0.06, and 0.12/day. Each result is divided by the same baseline integral.
When $\kappa=0$, density is independent of the signal even when the signal
changes. This is an executable limiting case, not biological validation.
Notebook 06 additionally shows that a constant one-way signal can be
absorbed into $\rho_{\mathrm{effective}}=\rho(1+\kappa b)$.
The same density trajectory can therefore leave the mechanism
undetermined within this model family.

\IfFileExists{generated/neuron_glioma_sensitivity.pdf}{
\begin{figure}[htbp]
\centering\includegraphics[width=0.62\textwidth]{generated/neuron_glioma_sensitivity.pdf}
\caption{Prescribed sensitivity in the synthetic model. Ratios concern
normalized-density integrals and are neither treatment effects nor
patient risk estimates. Coupling coefficients are varied, not inferred.}
\label{fig:sensitivity}
\end{figure}
}{\draftnotice{The sensitivity figure has not been generated.}}

\section{Clinical evidence and a proposed cognitive study}

Imaging and biological measurements require an explicit observation model
before they can calibrate these constructions. Structural MRI does not
directly identify immune composition or neuron--glioma synapses.
Functional connectivity is not interchangeable with synapse count.
The phenomenological index $a$ has no established mapping to EEG or MEG.

\begin{table}[htbp]
\centering\small
\caption{Proposed measurements for testing hypotheses. These are research
requirements, not observations collected for this perspective.}
\begin{tabularx}{\textwidth}{@{}>{\raggedright\arraybackslash}p{0.25\textwidth}
>{\raggedright\arraybackslash}X>{\raggedright\arraybackslash}X@{}}
\toprule
Measurement & Question it could address & Further requirement\\
\midrule
Reviewed longitudinal MRI & Geometry and change in defined compartments
& Treatment context, uncertainty, and an observation model\\
Spatial RNA/protein assays & Cell programs and niche composition
& Source validation, tissue sampling, subtype, and spatial correspondence\\
Physiological recordings & Aspects of activity and connectivity
& A justified relation to the modeled signal and confounding assessment\\
Functional perturbation & Whether a proposed interaction changes behavior
& Suitable controls and evidence beyond association\\
Clinical follow-up & Patient-relevant performance or benefit
& An eligible population, held-out evaluation, and an appropriate comparison\\
\bottomrule
\end{tabularx}
\end{table}

RANO 2.0 supplies clinical-trial response-assessment context
~\cite{wen2023}; a simulated density change is not a response category.
Clinical-record eligibility and source interpretation remain essential but
are not exercised here. No survival model, genotype inference, or
MRI--behavior cohort linkage is introduced.

\subsection{A proposed question about cognitive function}

\aicontribution
{Do preoperative network/dynamics features add information about later cognitive performance (notebooks 08/10)?}
{Proposed compact predictive AI and assistance with source linkage, feature definitions and leakage checks; neither predictor nor benefit is evaluated.}
{SSP follow-up span is a candidate primary endpoint; eligibility, window, feature budget and clinical review must be frozen prospectively.}
{Compare persistence, clinical baseline, anatomy and added network information using patient-level held-out evaluation and training-fold-only preprocessing.}
{Study design for a larger linked glioma cohort; the present behavioral example is not predictive validation.}

We propose asking whether preoperative network and dynamics information adds
predictive value beyond baseline cognition, clinical information and reviewed
anatomy. Candidate primary outcome is follow-up CANTAB SSP span, with baseline
SSP included as a predictor; RTI, RVP and SOC are secondary outcomes with their
own units. The endpoint and follow-up window require neuropsychological and
investigator review. The current behavioral availability audit is too small
for credible flexible model comparison and does not establish imaging linkage.
A larger linked glioma cohort and prospectively justified feature budget are
required. No molecular subtype is inferred from historical grade.

A persistence prediction retains baseline SSP. Subsequent comparisons add
available preoperative clinical variables, a prespecified anatomical feature,
then a few reviewed structural/functional-connectivity summaries or
independently fitted baseline dynamics parameters. A compact regularized
supervised model is the proposed predictive AI component. Keeping the same
model family across feature sets distinguishes added information from added
flexibility. Conventional connectivity analysis and The Virtual Brain
simulation remain distinct from AI. Earlier virtual-resection work showed
variable agreement with postoperative functional connectivity, rather than
established cognitive prediction~\cite{aerts2020}.

All predictors must be available at the preoperative prediction time.
Actual postoperative treatment, resection cavities and follow-up connectivity
cannot enter that forecast. A separate retrospective analysis could relate
their changes to cognitive changes, explicitly as association. Assessment
timing, treatment, seizures, medication, baseline function, education, subtype,
motion and hemodynamic-response handling need explicit documentation.

Use common eligible held-out patients for model comparisons. Scaling,
imputation, data-driven selection and tuning belong inside training folds;
all visits from one person remain together. Report held-out absolute errors
in span units, paired error differences against baselines, prediction failures,
missingness and uncertainty. Validation folds and repeats do not increase the
number of independent patients. Sample size and model complexity should be
chosen prospectively through justified simulation; this perspective supplies
no assumed effect size or achieved predictive result. Blank local study
schemas and the detailed protocol accompany the code.

An EEG/MEG/ECoG extension could test whether measured wave organization adds
information, where appropriate recordings exist. BOLD and fitted neural-mass
parameters are not direct LFP or field measurements. The present CANTAB
example motivates behavioral endpoints for analog-cognition hypotheses;
it does not discriminate that account from other network explanations.
Predictive accuracy would not alone establish mechanism or patient benefit.

Testing cytoelectric coupling is a separate experimental problem. Spatial
extracellular measurements would need concurrent molecular imaging and
controlled perturbation, specifying field gradient, frequency, orientation,
exposure and cell type. Controls would distinguish direct field effects from
altered spiking, synaptic/ionic signaling, temperature and viability.
Correlation alone would not establish the proposed intracellular mechanism.

\section{Methods and evidence boundaries}

\subsection{Shared code and reproducible outputs}

The code uses Python and a uv lockfile. Scientific functions live in
\texttt{src/glioma\_ai/}; notebooks call the same functions as the command-line
interface. The default examples run on CPU without patient files, weights,
API credentials, or a model endpoint. Each run records units, parameters,
software versions, source hashes, and generated-file hashes.

The intended primary workstation is a 128\,GB MacBook Pro with M4 Max.
Its capacity is a resource description, not a performance measurement.
Optional PyTorch device checks are separate from the NumPy examples. A
private-repository source archive provides a Colab path without embedding
GitHub credentials. A platform is called validated only after an executed
check is recorded; configuration alone is not a passed run.

\subsection{Roles and acceptance evidence}

\aicontribution
{What makes a contribution inspectable (notebook 04)?}
{Research assistance recording actual development and connecting the potassium record to checks.}
{The investigator requested explicit roles; human scientific approval remains a separate decision.}
{Record completeness and independent analytic/balance checks; neither supplies review by itself.}
{Retrospective development record; no comparative assistance experiment completed.}


\begin{table}[htbp]
\centering\small
\caption{Roles kept separate in the workflow. Initial synthetic examples
do not run a trained imaging model or a language model.}
\begin{tabularx}{\textwidth}{@{}p{0.22\textwidth}XX@{}}
\toprule
Role & Artifact & Relevant check\\
\midrule
Imaging model & Compartment label map & Label definitions, geometry, source
images, and appropriate reference evaluation\\
Proposed cognitive predictor & Follow-up spatial-span forecast & Held-out
patient errors against cognitive/anatomical baselines, uncertainty and failures\\
Research assistant & Proposed task, code, or prose & Source verification,
independent expected values, and researcher review\\
Numerical solver & Prescribed density trajectory & Discretization,
stability, conservation where applicable, and parameter sensitivity\\
\bottomrule
\end{tabularx}
\end{table}

The assistance record retains the task, permitted inputs, visible tool
information, contribution, artifact reference, checks, decision, and unresolved
questions. Human decisions name a reviewer only when review has occurred.
A record validator checks completeness, not truth. Retrospective development
summaries are labeled as summaries; they are not experimental transcripts.

The three unmeasured endpoints are productivity, scientific quality, and
clinical utility. Prospective timing and correction effort could address
productivity. Independent artifact review could assess scientific quality.
Patient benefit would require a patient-relevant endpoint and an appropriate
comparison. A correct figure does not by itself establish improved care.

\subsection{Prospective evaluation of research assistance}
\label{sec:ai-pilot}
The proposed protocol in \path{docs/ai-benefit-evaluation.md} separates
productivity, scientific quality and clinical relevance. A feasibility pilot
would involve two investigators, each completing six matched task pairs:
24 planned sessions, not collected observations or a powered efficacy study.
One distinct version of each pair permits a specified generative assistant;
the other uses the same code, editor, documentation and sources with
generative assistance disabled. Version assignment and condition order would
be counterbalanced and randomized before work begins. New tasks, practice
sessions and recorded prior expertise reduce avoidable learning effects;
carryover remains a limitation.

Task families cover imaging measurement, potassium, field units, growth
controls, fabricated clinical eligibility and evidence communication.
Freeze the task packets, independent acceptance key, tool information and
protocol before collection. Record only identifiers and settings actually
visible, retaining prompts and corrections locally where permitted.

The primary measure is quality-qualified completion time with a 60-minute
cap, including assistant latency, verification and self-correction. An
unaccepted or unfinished submission receives the cap. After submission,
independent reviewers blinded to condition would rate units/geometry,
correctness, controls/reproducibility, provenance and evidence boundaries
on 0--2 scales. Acceptance requires mandatory checks, at least 8/10 and
no critical error. Preserve independent ratings and disagreement resolution.

Report every assisted-minus-unassisted paired difference, within-investigator
medians, acceptance and critical-error counts. Do not exclude failed attempts
or treat sessions as independent participants. Quality scores and written
clinical-relevance judgments are separate outcomes; no patient outcome is
measured. Blank recording forms are supplied, with no assignments, timings
or reviews. Participants, task packets, reviewers and applicable institutional
requirements remain to be finalized before any pilot begins.

\section{Discussion and research agenda}

The contribution of these examples is a chain another researcher can
inspect: question, specification, contribution, decision, check, and output.
The default program creates deterministic constructions small enough for
teaching. The package and notebook interface reduce ambiguity about which
code produced a number, while the provenance record identifies the source
and environment of a run.

The potassium case connects an actual retrospective contribution with code,
independent checks and interpretation. It does not resolve the missing
interaction-level history. Future work should retain prospective inputs,
proposals, decisions and resulting artifacts as they occur, including the
possibility that no correction is necessary. The matched-task protocol makes
comparative benefit testable without inventing failures or approvals.

The source-verified CANTAB example adds empirical task performance to this
chain. Verification narrowed what the prepared table can support: extra local
digits cannot establish a change finer than the original release. Native
domain scales, assessment timing and missingness remain visible. These
observations anchor the clinical question without identifying a surgical or
growth effect. A linked network/dynamics study would test added predictive
information beyond baseline cognition and anatomy; the present availability
audit is not sufficient predictive validation.

For the conference, the glioma--brain interface organizes the narrative.
A niche example establishes why composition matters; potassium clearance and slow growth
feedback supply separate worked models. Imaging and clinical questions
anchor interpretation, while software mechanics and record-counting exercises
move to backup material. A proposed collaboration would combine reviewed
imaging with appropriate biological measurements and independent model review;
no such study is reported as agreed or completed.

The limitations are substantial. Synthetic shapes omit biological,
acquisition, and treatment variability. The density equation has prescribed
coefficients, no treatment term and no coupling to the surrounding brain.
Forecasting growth in the sense of weather prediction would further require
data assimilation from serial scans, uncertainty ensembles and validation
against held-out follow-up examinations. The niche construction aggregates
cell states and assumes reversibility; the neural extension has no explicit
neurons, synapses, or immune dynamics. The potassium construction has no
spatial buffering or explicit cells; the potential example supplies no
intracellular field. Neither predicts cognitive performance. The analog
cognition and cytoelectric interpretations remain hypotheses. No biological calibration or MRI
forward model is supplied. The source-presence test is narrower
than a clinical extractor. Optional public workflows depend on separately
prepared data. None of these artifacts measures the comparative benefit
of AI assistance.

The proposed pilot compares distinct matched tasks with independent quality
criteria and recorded time. Its small within-investigator design cannot
establish general superiority. Clinical utility would require an eligible
population, patient-relevant endpoints and an appropriate comparison in a
separate clinical study.

\section{Conclusion}
The glioma--brain interface connects published anatomical imaging with
questions about cell states, astrocytes, neural dynamics and cognitive function.
The executable models keep their units, timescales and observation boundaries
explicit. Experimental mechanisms, patient associations and emerging
field/cytoskeleton theories motivate different tests. AI assistance supports
inspectable research artifacts; comparative research benefit and clinical
utility remain unmeasured. The source-verified CANTAB example anchors cognition
in measured task performance. A proposed compact AI predictor would test
whether preoperative network information adds value beyond baseline cognition
and anatomy in a larger linked glioma cohort. Electrophysiological and
cellular hypotheses require their own measurements and controlled tests.

\section*{Declarations}

\paragraph{AI assistance.}
A Codex desktop assistant helped design this repository and draft the
adapted code, examples, notebooks, tests, documentation, and manuscript.
The development notes and worked potassium case are retrospective summaries.
The current framing and pilot design have a separate contribution record.
The empirical CANTAB verification, implementation and prediction-study design
also have a separate retrospective record, with endpoint and model choices
pending investigator and neuropsychological review.
The saved records do not identify an exact backend snapshot or sampling settings.
Rebuilding the examples does not call a language model. Human scientific
review remains pending.

\paragraph{Code and data.}
The development repository is
\url{https://github.com/arvidl/AI-assisted-glioma-research-dev}, initially
private. Selected helpers were adapted under the retained MIT license
from the predecessor repository at commit
\path{f8a37993d11f595135f407638750442041422391}.
No acquired patient files, original presentations, agenda, or invitation
are included. Selected published public-data illustrations are reused locally
under the article's CC BY 4.0 license, with attribution and a checked source
manifest. They are distinct from generated synthetic results. Public data
have their own release terms.
The CANTAB source revisions are pinned separately. Acquired behavioral tables,
local case identifiers, exact dates, individual scores and the generated
empirical case figure remain local and excluded from Git/source archives.
The local illustrated PDFs are separate from the source-only build variant.

\paragraph{Authorship, funding, and competing interests.}
The author and University of Bergen affiliation follow the investigator's
instructions. Collaborator contributions, funding, competing interests,
and target journal remain to be agreed and completed. No co-author
approval or clinical evaluation is asserted.

\appendix
\section{Reproduction and educational notebooks}

Install the locked environment and run:
\begin{verbatim}
uv sync --locked --group notebooks
make published-figures HPGS_REPO=/path/to/author-checkout
make check
make notebooks
make paper
make talk
\end{verbatim}
The published-figure import is optional. To prepare the empirical cognition
branch, explicitly run \texttt{make cantab-fetch}, create ignored
\path{local/cantab-case.json} with an explicit source ID and prepared session-table
path, then run \texttt{make cantab} before the notebooks and PDF builds.
The preparation instructions retain source revisions, hashes and license.
Without those inputs, the builds show notices rather than an empirical figure.
Notebook 01 covers imaging measurements and coordinate geometry; 02 supplies
the density baseline; 04 covers assistance records; 05 covers a
niche-dependent state transition; 06 is the slow growth-feedback experiment;
07 covers astrocytic potassium clearance; 08 covers potential, field and
cognitive-evidence boundaries; 09 supplies the published imaging bridge and
optional attributed illustrations; 11 covers network disconnection and model
functional connectivity on a synthetic connectome; 12 applies that model
descriptively to one patient's openly licensed (CC0) pre- and postoperative
connectomes, without reviewed upstream quality control. The main learning route is
04, 09, 01, 05, 07, 08, 02, 11, 12, 06, 10: first inspect the AI roles and worked record,
then follow imaging into biological modeling. Notebook 03 requires explicit
local public-data configuration. Notebook 00 provides the upload-based Colab path.
Notebook 10 checks the pinned behavioral source, local score agreement and
cohort availability; its empirical case needs the explicit make cantab preparation.
No acquisition or synthetic substitution occurs when local inputs are absent.
\aicontribution
{How can a private source archive reproduce the examples (notebook 00)?}
{Research assistance with bootstrap code and provenance documentation.}
{The investigator requested uv and Colab compatibility; permitted cloud inputs remain a decision.}
{Locked installation, source revision and local export reproduction.}
{One hosted Colab CPU session completed every cell before the notebooks were renumbered; no language model is called during rebuilding.}
Executed copies of notebooks 01--12 are tracked in \path{notebooks/executed/}
with their source licences; other generated results remain local. The
validation document records actual execution and skips.

\bibliographystyle{unsrtnat}
\bibliography{references}
\end{document}
